% ============================================================ Capítulo 1
\chapter{Introducción y contexto empresarial}\label{cap:intro}

\section{La empresa}
Fábrica Nacional de Autopartes ---Fanalca--- es un grupo industrial colombiano con sede principal en la
Calle~13 \#~31A-80, zona industrial de Acopi, Yumbo (Valle del Cauca). Su origen se remonta a la Fábrica
Nacional de Carrocerías, constituida el 31 de mayo de 1958 por las familias Losada y Herrera con doce
colaboradores~\cite{fanalca-quienes}. A lo largo de más de seis décadas la organización se diversificó hacia
las autopartes, la tubería de acero, el ensamble y la comercialización de motocicletas y automóviles Honda,
los equipos para la recolección de residuos sólidos y la infraestructura vial, con aliados comerciales en más
de dieciséis países~\cite{fanalca-quienes}. Fuentes de prensa regional describen el negocio Honda como el de
mayor participación en los ingresos de la organización~\cite{larepublica-fanalca}.

\begin{table}[H]\centering\small
\caption{Unidades de negocio publicadas en el sitio oficial}\label{tab:unidades}
\begin{tabularx}{\linewidth}{@{}lX@{}}\toprule
Unidad & Descripción resumida (fuente: páginas de negocio de \url{https://fanalca.com/en/business/})\\\midrule
Honda Motocicletas & Única planta nacional en el mundo con licencia para fabricar y ensamblar motocicletas
Honda; portafolio de 100 a 250\,cc; red de 450 puntos de venta en 32 departamentos.\\
Honda Autos & Importación y venta de vehículos Honda; ``Garantía Total Verdadera'' de 2 años sin límite de
kilometraje; 26 vitrinas.\\
Tubos y perfiles & Tubería de acero soldada y perfiles abiertos para construcción; plantas en Acopi-Yumbo y
Guachené (Cauca).\\
Autopartes & Partes metálicas para automóviles y vehículos comerciales; plantas en Yumbo, Bogotá y Medellín.\\
Ambiental & Cajas compactadoras y contenedores de acero desde 1981; presencia en 29 departamentos.\\
Fanalvías & Barreras metálicas, amortiguadores de impacto y terminales para seguridad vial.\\\bottomrule
\end{tabularx}
\end{table}

\section{El problema}
El Taller~1 plantea la ``necesidad de un canal de comunicación automatizado y preciso para la empresa''
(Taller~1, p.~2, §5). Un visitante que contacta por primera vez busca información dispersa en varias páginas
(sedes, unidades de negocio, historia, contacto, procesos básicos). Un asistente generativo sin control podría
responder con datos plausibles pero inexistentes ---por ejemplo, horarios de atención que el sitio no
publica---, lo cual compromete la confianza en la marca. El Módulo~1 aborda la capa que precede al
\emph{chatbot}: la base de conocimiento y un sistema Q\&A cuyas respuestas se limitan a lo que la empresa ha
publicado.

\section{Enfoque de la solución}
La solución se organiza en cuatro bloques (Figura~\ref{fig:arquitectura}):
\begin{enumerate}
  \item \textbf{Extracción responsable}: rastreador que respeta \texttt{robots.txt}, pausas, dominios
        permitidos y deja trazabilidad completa de cada captura.
  \item \textbf{Arquitectura Medallion}: Bronze (HTML original), Silver (texto limpio y deduplicado) y Gold
        (Markdown con metadatos, fragmentos citables y archivo consolidado).
  \item \textbf{Tareas LLM con LangChain}: Resumen, FAQ y Q\&A con tres versiones de prompt, verificación de
        citas y bloqueo de inyección de instrucciones.
  \item \textbf{Interfaces y gobierno}: Streamlit para la demostración, API REST, CLI, servidor MCP para
        agentes de código, auditoría en SQLite y un proceso de desarrollo asistido por OpenCode y Claude Code.
\end{enumerate}

\begin{nota}[Corpus incluido en el repositorio]
El entorno en el que se construyó el proyecto no tenía acceso de red al dominio \texttt{fanalca.com} desde la
línea de comandos. Por ello el repositorio incluye un \textbf{corpus semilla}: paráfrasis en español de diez
páginas públicas consultadas el 4 de noviembre de 2026 mediante una herramienta de lectura web, cada una con su
URL de origen. La semilla permite ejecutar la aplicación y las pruebas sin red. La base exigida por el taller
(``proveniente del web scraping'') se genera en el computador local con \cmd{uv run fabrica pipeline}
(capítulo~\ref{cap:procesodatos}); cada fragmento conserva el campo \texttt{origen} para auditar
qué información respalda cada respuesta.
\end{nota}

\section{Organización del documento}
Los capítulos~\ref{cap:intro} a~\ref{cap:requisitos} presentan el análisis documental y los requisitos; los
capítulos~\ref{cap:arquitectura} a~\ref{cap:pila} la arquitectura y la pila tecnológica; los
capítulos~\ref{cap:windows} a~\ref{cap:git} la instalación y configuración del entorno agéntico; los
capítulos~\ref{cap:scraping} a~\ref{cap:frontend} la implementación; los capítulos~\ref{cap:agentes} a
\ref{cap:problemas} la guía operativa, la calidad, la seguridad y la operación; el capítulo~\ref{cap:validacion}
contrasta el resultado con la rúbrica e incluye el informe exigido por el taller. Los capítulos~\ref{cap:simulador}
y~\ref{cap:anatomia} documentan la versión~2.0: el simulador 3D de ensamblajes y la biblioteca de máquinas.

% ============================================================ Capítulo 2
\chapter{Objetivos y alcance}\label{cap:objetivos}

\section{Objetivo general}
Diseñar y construir el núcleo de conocimiento para un futuro asistente virtual de Fanalca: extraer información
pública, procesarla y crear un sistema básico de Preguntas y Respuestas con técnicas de Prompt Engineering
(objetivo literal del Taller~1, p.~1).

\section{Objetivos específicos}
\begin{enumerate}[label=OE\arabic*.]
  \item Identificar las fuentes públicas de la empresa y definir el alcance de las preguntas de un cliente en
        su primer contacto.
  \item Extraer mediante \emph{web scraping} el contenido textual relevante del sitio oficial.
  \item Limpiar, normalizar y segmentar el texto en fragmentos semánticamente coherentes.
  \item Consolidar el conocimiento en un archivo de texto limpio e inyectarlo en el prompt de sistema.
  \item Implementar Resumen, FAQ y Q\&A con LangChain y prompts que obliguen a basarse en el contexto.
  \item Documentar la experimentación con prompts (zero-shot, formato, instrucciones anti-alucinación).
  \item Construir una interfaz Streamlit que muestre claramente las tres funcionalidades.
  \item Evaluar con al menos veinte preguntas distintas y documentar resultados y limitaciones.
\end{enumerate}

\section{Alcance}
\begin{table}[H]\centering\small
\caption{Delimitación del alcance}\label{tab:alcance}
\begin{tabularx}{\linewidth}{@{}>{\bfseries}l X@{}}\toprule
Documentado (Módulo 1) & Scraping del sitio oficial; limpieza; \emph{chunking}; archivo consolidado; contexto
completo en el prompt de sistema; Resumen, FAQ y Q\&A con LangChain; Streamlit; 24 preguntas de prueba; informe.\\
Decisiones de diseño & Arquitectura Medallion en archivos; proveedores LLM intercambiables; prompts versionados
en YAML; salvaguardas posteriores a la generación; banco de evaluación con métricas automáticas.\\
Extensiones propuestas & Selección de contexto BM25; API FastAPI; servidor MCP; auditoría SQLite con Alembic;
línea base extractiva; contenedores Docker.\\
Fuera de alcance & Embeddings y base vectorial (Módulo~2); memoria conversacional y enrutamiento (Módulo~2);
canales WhatsApp/Telegram (Fase~3); scraping de redes sociales.\\\bottomrule
\end{tabularx}
\end{table}

\section{Alcance del Q\&A}
Según la instrucción ``Definición del Alcance'' (Taller~1, p.~1, §1), el asistente atiende preguntas sobre:
sedes y plantas, unidades de negocio y productos, historia, información de contacto, procesos básicos
(formulario de contacto, vacantes, garantía) y sostenibilidad. Las preguntas cuya respuesta no está publicada
---por ejemplo, horarios o precios--- forman parte del alcance \emph{de prueba}: el comportamiento correcto es el
rechazo explícito.

% ============================================================ Capítulo 3
\chapter{Análisis de los documentos}\label{cap:documentos}

\section{Documentos examinados}
\begin{table}[H]\centering\small
\caption{Inventario documental}\label{tab:documentos}
\begin{tabularx}{\linewidth}{@{}l c X@{}}\toprule
Documento & Págs. & Contenido relevante\\\midrule
Taller 1 --- Aplicación de Técnicas Avanzadas de IA (dos copias) & 1--3 & Objetivo, cinco instrucciones,
rúbrica ajustada y criterios de calificación\\
Clase 1 --- Introducción a los LLMs & 1--4 & Herramientas del curso; proyecto integrador en tres fases\\
Clase 2 --- Modelos cerrados vs. abiertos y uv & 1--4 & uv; arquitectura Medallion; scraping con criterio\\
Clase 3 --- Consolas agénticas (OpenCode) (dos copias) & 1--5 & OpenCode; MCP; skills; SkillSpector; Engram;
OpenSpec; GGA\\
Prompt maestro del proyecto & --- & Requerimientos de ingeniería del entregable (este manual)\\\bottomrule
\end{tabularx}
\end{table}
Los PDF provienen de la impresión de páginas de la plataforma UAO Virtual; su texto no es extraíble porque las
fuentes usan una codificación propia. El análisis se realizó sobre las páginas rasterizadas, verificando cada
cita contra la imagen.

\section{Requisitos extraídos literalmente}
\begin{itemize}
  \item ``Utilizando librerías de Python como requests y BeautifulSoup, Selenium, los estudiantes deberán hacer
        web scraping del sitio web de su empresa'' (Taller~1, p.~1, §2).
  \item ``El texto extraído debe ser limpiado (eliminar HTML, etc.) y dividido en fragmentos más pequeños y
        semánticamente coherentes (chunks)'' (p.~1, §2).
  \item ``Diseñarán un prompt robusto que instruya al LLM para que responda la pregunta del usuario basándose
        únicamente en el contexto proporcionado'' (p.~1, §3).
  \item ``Tomen todo el texto limpio que extrajeron en el `Punto 2' [\ldots] y consolídenlo en el prompt de
        sistema'' (p.~1, §3).
  \item ``Formulando al menos 20 preguntas distintas (relacionadas con el alcance definido en el paso 1)''
        (p.~2, §5).
  \item ``Implementación funcional de las tres tareas solicitadas (Resumen, FAQ, Q\&A) usando LangChain''
        (p.~2, rúbrica).
  \item ``Se respeta robots.txt (lo que está en Disallow no se toca), se identifica el User-Agent, se pone
        timeout y un sleep entre peticiones'' (Clase~2, p.~3).
  \item ``nada de pip ni requirements.txt'' (Clase~2, p.~2).
\end{itemize}

\section{Actores, datos y entregables}
\textbf{Actores}: cliente o visitante (consulta), equipo de desarrollo de cuatro personas (construcción),
docente evaluador (revisión y sustentación). \textbf{Sistemas externos}: sitio \texttt{fanalca.com}, servidor
Ollama o API de LLM, GitHub. \textbf{Datos}: secciones ``Quiénes somos'', negocios, sostenibilidad, noticias y
contacto. \textbf{Entregables}: repositorio con código y base de conocimiento; aplicativo; interfaz; pruebas
con 20 o más preguntas; informe con cinco secciones; sustentación en vivo de 15 minutos sin PowerPoint.

\section{Restricciones}
Académicas (grupos de cuatro; apertura 3 de abril y cierre 2 de mayo de 2026, con la copia del 24/10/2026
marcando la entrega como vencida), técnicas (LangChain recomendado, Ollama o API privada, uv obligatorio) y
éticas o legales (datos públicos, \texttt{robots.txt}, Ley~1581 de 2012 de protección de datos
personales~\cite{ley1581}).

\section{Información ausente, ambigua o contradictoria}
\begin{longtable}{@{}P{0.6cm}P{7.2cm}P{7.4cm}@{}}
\caption{Hallazgos del análisis documental}\label{tab:ambiguedades}\\\toprule
\# & Hallazgo & Tratamiento adoptado\\\midrule\endfirsthead
\toprule \# & Hallazgo & Tratamiento adoptado\\\midrule\endhead
A1 & La sección 4 del taller queda cortada entre páginas; solo es legible ``\ldots muestre la respuesta
generada por el sistema''. & Interfaz que recibe la pregunta y muestra la respuesta, coherente con el título
``Desarrollo de una Interfaz de Prueba''.\\
A2 & Dos esquemas de evaluación: rúbrica ajustada (15/25/25/15/10/10~=~100\,\%) y tabla de criterios
(10/20/20/20/10/10 visibles, 90\,\%, tabla truncada). & Se cubren ambos; la validación del
capítulo~\ref{cap:validacion} sigue la rúbrica ajustada.\\
A3 & La rúbrica exige Resumen y FAQ; las instrucciones solo detallan Q\&A. & Se implementan las tres tareas.\\
A4 & El informe pide el modelo de \emph{embeddings} y la base vectorial ``a utilizar en el módulo 2'', mientras el
Módulo 1 consolida todo en el prompt. & Contexto completo en el Módulo 1; propuesta formal para el Módulo 2 en
OpenSpec.\\
A5 & La Clase 2 nombra las capas \texttt{data/raw}, \texttt{data/processed}, \texttt{data/gold}; el prompt
maestro usa \texttt{bronze/silver/gold}. & Se adopta Bronze/Silver/Gold y se documenta la equivalencia
(ADR-002).\\
A6 & Asignación: ``Top 50'' asignado (taller) frente a elección entre ``16 empresas del ranking'' (Clase 1). &
Sin impacto técnico.\\
A7 & Razón social ``S.A.S.'' (solicitud) frente a ``Fanalca S.A.'' (pie del sitio). & Se usa ``Fábrica
Nacional de Autopartes (Fanalca)''.\\
A8 & El alcance sugiere ``horarios'', pero el sitio no los publica. & La pregunta P20 mide el rechazo
correcto.\\
A9 & El taller agrupa ``GPT-5, Gemini 2.5 Pro'' bajo ``API de OpenAI''. & Se interpreta como API privada; se
soportan OpenAI y Anthropic.\\\bottomrule
\end{longtable}
